\subsection{Fine-Grained Do-No-Harm}
\label{sec:evaluation:fine-donoharm}

The do-no-harm property reported in
\S\ref{sec:evaluation:topology} and
\S\ref{sec:evaluation:realtopology} (\GreenFaaS{} = FIFO byte-for-byte
in single-region or zero-shiftable settings) admits a fair concern:
when the workload offers \emph{no} shifting opportunity, every
reasonable scheduler is trivially identical to FIFO. A more
substantive test is whether \GreenFaaS{} degrades \emph{gracefully}
across the regime of very small but nonzero shifting opportunity ---
where the scheduler actually has decisions to make but the carbon
upside is marginal, and a naively eager scheduler could pay
non-trivial idle-energy cost for negligible savings.

\begin{table}[h]
\centering\small
\caption{Fine-grained shiftable-fraction sweep. Real LWA carbon,
5 regions, 6h, $\sim$43k invocations. ``Gap'' is carbon reduction vs
FIFO (positive = better than FIFO).}
\label{tab:fine-donoharm}
\begin{tabular}{rrrrr}
\toprule
Shiftable frac. & FIFO (g) & GreenFaaS gap & v1 gap & Spatial gap \\
\midrule
0.00          & 6.20 & +0.00\%  & +0.00\%  & +0.00\%  \\
0.01          & 6.20 & +0.00\%  & +0.00\%  & +0.00\%  \\
0.02          & 6.20 & +0.00\%  & +0.00\%  & +0.00\%  \\
0.05          & 6.20 & +3.64\%  & +4.14\%  & +4.14\%  \\
0.10          & 6.20 & +3.64\%  & +4.14\%  & +4.14\%  \\
0.25          & 6.20 & +21.3\%  & +22.0\%  & +22.0\%  \\
\bottomrule
\end{tabular}
\end{table}

The findings refine the do-no-harm claim into something more
substantive. At very low shiftable fractions (0, 1\%, 2\% of the
workload), \GreenFaaS{}, GreenFaaS-v1, and Spatial all match FIFO
\emph{byte-for-byte}. This is not the trivial property
``no-shifting-implies-FIFO'' but the stronger property that
\GreenFaaS{} correctly identifies negligible-opportunity regimes and
declines to act. As the shiftable fraction grows past a small
threshold (around 5\% here), the carbon-aware schedulers begin
diverging from FIFO smoothly, with \GreenFaaS{} producing slightly
smaller savings than Spatial or v1 in this region --- the
\S\ref{sec:tradeoff:idle} correction is conservative when in doubt,
which is the desired behaviour for a deployment-ready scheduler.

The regime where \GreenFaaS{}'s conservatism becomes an advantage
appears at higher shiftable fractions where carbon variability is
high and idle-energy cost is non-trivial; the
\S\ref{sec:evaluation:variability} variability sweep (and the v1
collapse documented there) is the empirical evidence for this.

The interpretation: do-no-harm is not just ``correctly idle
when there are no decisions to make,'' but ``correctly idle across a
neighbourhood of low-opportunity workloads,'' which is the practical
shape of the property that matters for deployment.
